\chapter{Surgical Abortion}

\section*{Introduction \& Clinical Indications}
Surgical abortion is a medical procedure designed to terminate a pregnancy by removing the contents of the uterus. This procedure can be performed in the first or second trimester, with techniques varying based on gestational age. In the first trimester, vacuum aspiration, also known as suction curettage, is the most common method utilized, while in the second trimester, dilation and evacuation (D\&E) may be employed. Surgical abortion is typically conducted in an outpatient setting, such as a clinic or operating room, and is regarded as one of the safest surgical interventions when performed by a qualified healthcare provider. The primary aim of this procedure is to provide a safe option for individuals seeking to end an unwanted pregnancy, as well as to manage complications arising from early pregnancy, such as incomplete miscarriage or severe fetal anomalies.

\begin{itemize}
  \item Termination of pregnancy.
  \item Vacuum aspiration or suction curettage.
  \item Dilation and curettage (D\&C).
  \item Dilation and evacuation (D\&E).
  \item Aspiration abortion.
  \item Elective abortion.
\end{itemize}

The fundamental principle of surgical abortion is to achieve complete evacuation of the uterine contents while minimizing trauma to the uterine lining and surrounding structures. In the first trimester, the procedure typically involves the use of a suction cannula connected to a vacuum source, which gently aspirates the pregnancy tissue from the uterine cavity. Following suction, a curette may be employed to ensure that the uterine cavity is clear of any residual tissue. 

In the second trimester, the procedure necessitates more extensive cervical dilation, allowing for the removal of larger volumes of tissue using suction combined with surgical instruments such as forceps. The overarching goals of surgical abortion include ensuring complete evacuation of the uterus, reducing the risk of complications such as bleeding and infection, and facilitating a rapid recovery for the patient.

The practice of abortion has a long history, with various methods employed across cultures and eras. The modern surgical techniques began to take shape in the mid-twentieth century, particularly with the introduction of vacuum aspiration in the 1960s and 1970s. Harvey Karman and other pioneers played a significant role in advancing this technique, which greatly improved the safety and efficiency of first-trimester abortions compared to the sharp curettage methods used previously. 

Legal reforms in many countries during this period facilitated greater access to safe abortion services, significantly reducing morbidity and mortality associated with unsafe procedures. Over the decades, the integration of ultrasound technology has enhanced gestational dating and procedural safety, while the advent of medical (medication) abortion has provided an alternative option for early pregnancy termination.

Surgical abortion is indicated in various clinical scenarios, including:

\begin{itemize}
  \item Termination of an unwanted pregnancy at the request of the patient.
  \item Management of a pregnancy where continuation poses a serious risk to the health of the pregnant individual.
  \item Management of an anomalous pregnancy that is incompatible with life or associated with severe fetal abnormalities.
  \item Evacuation of the uterus after a miscarriage that has not completely expelled (incomplete abortion).
  \item Evacuation of retained products of conception following a spontaneous or medical abortion.
  \item Management of a molar (hydatidiform) pregnancy in selected cases.
  \item Cases of severe maternal health conditions that necessitate immediate termination.
  \item Situations where the patient has a strong psychological or emotional need to terminate the pregnancy.
\end{itemize}

Surgical abortion is primarily a first-line procedure for terminating a pregnancy. When a patient requests termination, surgical aspiration is one of the two main options available, the other being medical (medication) abortion. The choice between these options is influenced by factors such as gestational age, patient preference, and clinical circumstances. In cases of incomplete miscarriage or retained products of conception, surgical evacuation serves as an essential primary treatment to prevent complications such as heavy bleeding and infection.

\section*{Relevant Surgical Anatomy}
Understanding the relevant surgical anatomy is crucial for the safe performance of surgical abortion. The uterus is a hollow muscular organ located in the pelvis, supported by ligaments and surrounded by the bladder anteriorly and the rectum posteriorly. The cervix, which is the lower part of the uterus, opens into the vagina and serves as the passage for instruments during the procedure. 

The arterial supply to the uterus is primarily via the uterine arteries, which branch from the internal iliac arteries, while venous drainage occurs through the uterine veins, which drain into the internal iliac veins. The uterine lymphatic drainage follows the arterial supply, with lymph nodes located in the pelvic region. 

Key surgical landmarks include:

\begin{itemize}
  \item McBurney's point: Although primarily associated with appendicitis, knowledge of this landmark can aid in differentiating pelvic pain.
  \item Calot's triangle: Important in understanding the anatomy surrounding the gallbladder, which may be relevant in cases of concurrent surgical procedures.
\end{itemize}

The surrounding structures, including the bladder and rectum, must be carefully considered to avoid injury during the procedure.

\section*{Preoperative Workup \& Preparation}
Prior to the procedure, a thorough preoperative workup is essential. This includes counseling the patient, obtaining informed consent, and confirming gestational age through ultrasound when appropriate. A pregnancy test and review of the patient's medical history, including any bleeding disorders, are conducted. 

Preparation steps include:

\begin{itemize}
  \item Ensuring the patient is NPO (nothing by mouth) for a few hours if sedation or anesthesia is planned.
  \item Administering medication to soften the cervix in certain cases.
  \item Arranging for a support person to accompany the patient and for transportation home post-procedure.
  \item Providing contraceptive counseling, as fertility resumes quickly after the procedure.
\end{itemize}

Absolute contraindications to surgical abortion include:

\begin{itemize}
  \item An untreated severe bleeding disorder.
  \item The presence of an active pelvic infection, which must be treated prior to the procedure.
\end{itemize}

Relative contraindications and precautions include:

\begin{itemize}
  \item A known pregnancy of uncertain location, necessitating careful evaluation to rule out ectopic pregnancy.
  \item Uncertainty regarding gestational age or pregnancy location, requiring additional workup.
  \item The need to empty the bladder before the procedure.
  \item Availability of skilled staff to manage potential complications such as heavy bleeding or uterine perforation.
  \item Ensuring the patient understands the procedure and provides informed consent, with psychological support available if desired.
\end{itemize}

\section*{Surgical Technique \& Procedure}
The surgical abortion procedure is typically performed in a clinic or operating room setting, often utilizing local anesthesia combined with sedation. The patient is positioned in the lithotomy position on an examination table, and the cervix is visualized using a speculum.

The key steps in the procedure include:

\begin{itemize}
  \item Confirming gestational age through history, examination, and ultrasound.
  \item Cleaning the cervix with an antiseptic solution, followed by the administration of a local anesthetic around the cervix.
  \item Gradually dilating the cervix using dilators or administering medication to soften the cervix.
  \item Inserting a suction cannula through the cervical opening into the uterus for first-trimester aspiration, applying suction to evacuate the pregnancy tissue.
  \item Using a curette to ensure the uterine cavity is smooth and free of residual tissue.
  \item In the second trimester, performing more extensive cervical dilation and utilizing suction along with forceps and curettes to remove larger volumes of tissue.
  \item Inspecting the evacuated tissue to confirm complete removal and monitoring the patient for bleeding before discharge.
\end{itemize}

The entire procedure typically lasts only a few minutes for first-trimester cases, with pain managed through analgesics and local anesthesia.

\section*{Complications \& Risks}
The potential complications and risks associated with surgical abortion can be categorized as follows:

General surgical complications:

\begin{itemize}
  \item Bleeding, typically minor but occasionally severe.
  \item Infection of the uterus or pelvis, mitigated by aseptic technique and prompt evacuation.
  \item Anesthesia-related complications, which are rare.
\end{itemize}

Procedure-specific risks:

\begin{itemize}
  \item Incomplete evacuation, necessitating a repeat procedure.
  \item Uterine perforation, an uncommon but significant complication that may require observation or surgical intervention.
  \item Injury to the cervix, including cervical tears.
  \item The small risk of a continuing pregnancy or pregnancy of unknown location not being terminated.
  \item Emotional responses experienced by some patients post-procedure.
\end{itemize}

\section*{Postoperative Care \& Recovery}
Postoperative recovery begins in the post-anesthesia care unit (PACU), where the patient is monitored for vital signs, bleeding, and pain control. The first 24-48 hours typically involve mild cramping and spotting, which are expected. 

During the first week, patients are advised to:

\begin{itemize}
  \item Gradually resume normal activities as tolerated.
  \item Avoid intercourse, tampons, and swimming for one to two weeks or until bleeding resolves.
\end{itemize}

Long-term recovery generally occurs within a few weeks, with most patients returning to their regular routines shortly after the procedure.

During the procedure, the surgeon documents the condition of the uterine cavity and the nature of the evacuated tissue. The pathology report on any resected tissues provides critical information regarding the type of pregnancy and any abnormalities present. This documentation is essential for ensuring comprehensive patient care and follow-up.

A normal postoperative course following surgical abortion includes the resolution of pain, a gradual decrease in bleeding, and the return of the patient's body to its non-pregnant state. Patients typically experience a return of normal menstruation within four to six weeks, with no signs of infection or other complications.

Postoperative care involves scheduling follow-up visits to monitor recovery and address any concerns. Key aspects of postoperative care include:

\begin{itemize}
  \item Follow-up visits within one to four weeks to confirm complete recovery.
  \item Suture or staple removal if applicable.
  \item Rehabilitation or physical therapy as needed.
  \item Education on warning signs that warrant contacting a healthcare provider, such as heavy bleeding, fever, severe pain, or foul-smelling discharge.
\end{itemize}

Patients are encouraged to seek care for any concerning symptoms and to discuss contraceptive options during follow-up visits.

\section*{Outcomes \& Prognosis}
Surgical abortion is a prevalent procedure worldwide, with estimates indicating that tens of millions of abortions occur annually. The incidence of surgical abortion peaks during the reproductive years, with the majority of procedures performed in the first trimester. The safety of surgical abortion has improved significantly, and in settings where it is legally performed under hygienic conditions, the rate of serious complications is very low. Variability in incidence across countries is influenced by legal restrictions, access to contraception, and the overall healthcare infrastructure.

Surgical abortion is highly effective, with success rates exceeding 98 percent for first-trimester procedures. The small number of failures typically result from incomplete evacuation, which can be addressed with a repeat procedure. Serious complications, such as uterine perforation or significant bleeding, occur in less than 1 percent of cases. Recovery is generally swift, and most patients experience no long-term adverse effects on fertility. The prognosis following surgical abortion is excellent, with the primary long-term considerations being effective contraceptive use to prevent unintended pregnancies and addressing any psychological responses through support and counseling as needed.

\section*{References}
\begin{enumerate}
  \item MedlinePlus. Abortion. U.S. National Library of Medicine; 2025.
  \item Current Medical Diagnosis and Treatment. McGraw-Hill; 2025.
  \item Williams Gynecology, 4th edition. McGraw-Hill; 2020.
  \item American College of Obstetricians and Gynecologists Practice Bulletin on Medication Abortion and Surgical Abortion; 2020.
\end{enumerate}

\clearpage
